\scanpage{42}
% L1, upper left
$\overline{\underline{\mathrm{XII}}}$

% L1–4
\uline{Remarque} \inserted{Dans 6.5.} Les hypothèses de quasi-compacité et quasi-séparation ne servent que pour \deleted{\unsure{faire}} prouver l'équivalence de la condition a) avec les conditions en apparence plus faibles b), b') etc.

% L5–10
L'équivalence de ces conditions entre elles, et les implications \deleted{de} a) \deleted{\unsure{implique}} vers b) \unsure{énoncées|s'énonce} dans 6.5. bis, restent valables sans cette hypothèse et méritent d'être énoncées une bonne fois \ldots. Dans 6.6. également, on n'a \unsure{d'ailleurs} besoin de la quasi-compacité et quasi-séparation \unsure{;|,}

% L11–15
de plus, la catégorie de tous les faisceaux d'ens. (resp. de Groupes de torsion \unsure{$\mathscr L$|$\mathcal L$} etc) peut ad libitum s'y remplacer par telle ou telle sous-catégorie --- manifestement on n'a pas \unsure{donné|d'ennui} l'énoncé \unsure{bouffe-tout} abstrait.

% L15–23
Dans le lemme 6.7. de descente également, on n'utilise pas la quasi-compacité et quasi-séparation, (sauf, \uline{soit-disant}, pour passer à la limite \deleted{\unsure{dans la}} \inserted{à la catégorie} des faisceaux de $\unsure{\mathbb Z|Z}/n\mathbb Z$-Modules, mais c'est sans doute inutile \deleted{\unsure{quitte à en être affirmé}} \inserted{--- et par l'argument à la main, utilisant \uline{\unsure{VIII|VII|XIV}} 5.2.,} \deleted{\unsure{--- ce qui c'est que cette vérification}} pour déduire que l'hom de changement de base en dim $n+1$ est encore injectif --- mais cela doit aussi sortir directement.

% archive sheet number: 41
